
déjà
prouvé, on en conclut que quel que soit @@M0@@,
@@M1@@
pour tout @@M2@@.

Par 2.5, @@M3@@ est un @@M4@@-faisceau. Soient @@M5@@ pour @@M6@@ et @@M7@@. Par 1.11, @@M8@@ et @@M9@@ est @@M10@@-régulier. Donc par 1.4 (ii), @@M11@@ pour tout @@M12@@. Donc si l'on définit des polynômes par @@M13@@ et @@M14@@, alors on a (ii) en vertu de 2.9.2.

Enfin, comme @@M15@@ est @@M16@@-régulier, @@M17@@ pour @@M18@@ par 1.4 (i)\,; d'où @@M19@@. Mais @@M20@@ par 1.8 (i). Donc si l'on définit @@M21@@, on a (iii).

\textbf{Remarque 2.10.} Pour @@M22@@ fixé, considérons une égalité de polynômes en @@M23@@ :
@@M24@@
Alors on trouve facilement la formule
@@M25@@
en raisonnant par récurrence sur @@M26@@ afin que @@M27@@ devienne @@M28@@, et en faisant @@M29@@.

De même, on trouve qu'il existe un polynôme à coefficients rationnels @@M30@@ qui est croissant en @@M31@@, tel que dans l'égalité suivante :
@@M32@@
